\chapter{紧性}

\section{紧致空间}


\begin{definition}[开覆盖]
    设 $X$ 是一个拓扑空间，$A \subset X$.$X$ 中开集的一个集合 $\mathcal{O}$ 称为 $A$ 的一个\textbf{开覆盖}，如果 $A$ 包含在 $\mathcal{O}$ 中开集的并集之内.
\end{definition}

\begin{definition}[子覆盖]
    如果 $\mathcal{O}$ 覆盖 $A$，$\mathcal{O}'$ 是 $\mathcal{O}$ 的一个子集合且也覆盖 $A$，那么 $\mathcal{O}'$ 称为 $\mathcal{O}$ 的一个\textbf{子覆盖}.
\end{definition}

\begin{definition}[紧致空间]
    一个拓扑空间 $X$ 称为\textbf{紧致的}或\textbf{紧的}，如果 $X$ 的每一个开覆盖都有一个有限子覆盖.
\end{definition}

\begin{instance}
    $\mathbb{E}^1$ 不是紧的.开覆盖 $\mathcal{O} = \{(-n, n) \mid n \in \mathbb{Z}^+\}$ 没有有限子覆盖.$\mathbb{E}^n$ 也不是紧的.
\end{instance}

\begin{instance}
    如果 $X = \{x_1, x_2, \dots, x_n\}$ 是一个有限集，那么 $X$ 上的任何拓扑都是紧的.
\end{instance}

\begin{instance}
    $\mathbb{R}_{fc}$ (有限补拓扑下的实数集) 是紧的.
    \begin{proof}
        假设 $\mathcal{O}$ 是 $\mathbb{R}_{fc}$ 的一个开覆盖.任选一个 $U \in \mathcal{O}$，那么 $\mathbb{R} - U$ 是有限集，记作 $\{r_1, r_2, \dots, r_m\}$.对每一个 $r_i$，存在 $U_i \in \mathcal{O}$ 使得 $r_i \in U_i$.那么 $\{U_1, U_2, \dots, U_m, U\}$ 就是 $\mathcal{O}$ 的一个有限子覆盖.
    \end{proof}
\end{instance}

\begin{definition}[紧致子集]
    设 $X$ 是一个拓扑空间，$A \subset X$.如果 $A$ 在子空间拓扑中是紧的，则称 $A$ 在 $X$ 中是\textbf{紧致的}，或称 $A$ 是 $X$ 的一个\textbf{紧致子集}，简称\textbf{紧子集}.
\end{definition}

\begin{instance}
    $\mathbb{R}_{fc}$ 中的任何子集都是紧的.
\end{instance}

\begin{instance}[Heine-Borel 定理]
    $A \subset \mathbb{E}^n$ 是紧的 $\iff$ $A$ 是有界闭集.
\end{instance}

\begin{instance}
    $S^{n-1} \subset \mathbb{E}^n$ 是紧的.
\end{instance}

\begin{lemma}
    设 $X$ 是一个拓扑空间，$A \subset X$.$A$ 在 $X$ 中是紧的 $\iff$ $A$ 的任意由 $X$ 中开集组成的开覆盖都有一个有限子覆盖.
\end{lemma}
\begin{proof}
    “$\Rightarrow$” 假设 $A$ 在 $X$ 中是紧的，$\mathcal{O}$ 是 $A$ 的一个由 $X$ 中开集组成的开覆盖.令 $\mathcal{O}' = \{U \cap A \mid U \in \mathcal{O}\}$，这是 $A$ 的一个由 $A$ 中开集（根据子空间拓扑的定义）组成的开覆盖.因为 $A$ 是紧的，所以 $\mathcal{O}'$ 有一个有限子覆盖 $\{U_1 \cap A, U_2 \cap A, \dots, U_n \cap A\}$.那么 $\{U_1, U_2, \dots, U_n\}$ 就是 $\mathcal{O}$ 的一个覆盖 $A$ 的有限子覆盖.

    “$\Leftarrow$” 假设 $A$ 的任意由 $X$ 中开集组成的开覆盖都有一个有限子覆盖.设 $\mathcal{O}'$ 是 $A$ 在子空间拓扑下的一个开覆盖.根据子空间拓扑的定义，对任意 $V \in \mathcal{O}'$，存在 $X$ 中的开集 $U_V$ 使得 $V = U_V \cap A$.那么集合 $\mathcal{O} = \{U_V \mid V \in \mathcal{O}'\}$ 是 $A$ 的一个由 $X$ 中开集组成的开覆盖.根据假设，$\mathcal{O}$ 有一个有限子覆盖 $\{U_{V_1}, \dots, U_{V_n}\}$.于是，$\{V_1, \dots, V_n\}$ 就是 $\mathcal{O}'$ 的一个有限子覆盖.因此 $A$ 是紧的.
\end{proof}


\section{紧致空间的性质}

\begin{proposition}
    紧致空间的连续像是紧的.即若 $f: X \to Y$ 是连续函数且 $X$ 是紧致空间，则 $f(X)$ 在 $Y$ 中是紧的.
\end{proposition}
\begin{proof}
    假设 $\mathcal{O}$ 是 $f(X)$ 在 $Y$ 中的一个开覆盖.那么对任意 $U \in \mathcal{O}$，$f^{-1}(U)$ 在 $X$ 中是开集，并且 $\{f^{-1}(U) \mid U \in \mathcal{O}\}$ 是 $X$ 的一个开覆盖.因为 $X$ 是紧的，所以存在 $U_1, U_2, \dots, U_n \in \mathcal{O}$ 使得 $\{f^{-1}(U_1), f^{-1}(U_2), \dots, f^{-1}(U_n)\}$ 覆盖 $X$，即 $X = \bigcup_{i=1}^n f^{-1}(U_i)$.
    所以 $f(X) = f\left(\bigcup_{i=1}^n f^{-1}(U_i)\right) = \bigcup_{i=1}^n f(f^{-1}(U_i)) \subset \bigcup_{i=1}^n U_i$.因此 $\mathcal{O}$ 有一个有限子覆盖.
\end{proof}

\begin{corollary}
    设 $X_1, X_2$ 是拓扑空间，且 $X_1 \cong X_2$那么 $X_1$ 是紧的 $\iff X_2$ 是紧的.
\end{corollary}

\begin{instance}
    $\mathbb{E}^1 \not\cong \mathbb{R}_{fc}$ (非紧 $\leftrightarrow$ 紧)，$\mathbb{E}^n \not\cong S^n$ (非紧 $\leftrightarrow$ 紧).
\end{instance}

\begin{theorem}[极值定理]
    如果 $X$ 是一个紧致空间，$f: X \to \mathbb{E}^1$ 是连续函数，那么 $f$ 在 $X$ 上取得最大值和最小值.
\end{theorem}
\begin{proof}
    因为 $X$ 是紧的且 $f: X \to \mathbb{E}^1$ 是连续的，所以 $f(X)$ 在 $\mathbb{E}^1$ 中是紧的.这意味着 $f(X)$ 在 $\mathbb{E}^1$ 中是有界闭集.因此 $f$ 取得最大值和最小值.
\end{proof}

\begin{proposition}
    紧致空间的闭子集是紧的.
\end{proposition}
\begin{proof}
    设 $X$ 是紧致空间，$A \subset X$ 是闭集.假设 $\mathcal{O}$ 是 $A$ 的一个由 $X$ 中开集组成的开覆盖.因为 $A$ 是闭集，所以 $X-A$ 是开集.令 $\mathcal{O}_1 = \mathcal{O} \cup \{X-A\}$.那么 $\mathcal{O}_1$ 是 $X$ 的一个开覆盖.根据 $X$ 的紧性，$\mathcal{O}_1$ 有一个有限子覆盖 $\mathcal{O}_2$.
    如果 $X-A \notin \mathcal{O}_2$，那么 $\mathcal{O}_2$ 是 $\mathcal{O}$ 的一个覆盖 $A$ 的有限子覆盖.
    如果 $X-A \in \mathcal{O}_2$，那么 $\mathcal{O}_2 - \{X-A\}$ 是 $\mathcal{O}$ 的一个覆盖 $A$ 的有限子覆盖.
\end{proof}

\begin{proposition}
    Hausdorff 空间中的紧子集是闭的.
\end{proposition}
\begin{proof}
    设 $X$ 是 Hausdorff 空间，$A \subset X$ 是紧子集.我们证明 $X-A$ 是开集.固定一点 $x \in X-A$.对任意 $y \in A$，$x \neq y$.因为 $X$ 是 Hausdorff 的，所以存在 $x$ 的邻域 $U_y$ 和 $y$ 的邻域 $V_y$ 使得 $U_y \cap V_y = \emptyset$.
    集合 $\{V_y \mid y \in A\}$ 是 $A$ 的一个开覆盖.因为 $A$ 是紧的，所以存在有限个点 $y_1, \dots, y_n \in A$ 使得 $\{V_{y_1}, \dots, V_{y_n}\}$ 覆盖 $A$，即 $A \subset \bigcup_{i=1}^n V_{y_i}$.
    令 $U = \bigcap_{i=1}^n U_{y_i}$.$U$ 是 $x$ 的一个邻域（有限个开集之交）.
    我们有 $U \cap \left(\bigcup_{i=1}^n V_{y_i}\right) = \bigcup_{i=1}^n (U \cap V_{y_i})$.因为 $U \subset U_{y_i}$，所以 $U \cap V_{y_i} \subset U_{y_i} \cap V_{y_i} = \emptyset$.
    因此 $U \cap A = \emptyset$，这意味着 $U \subset X-A$.这就证明了 $X-A$ 是开集，所以 $A$ 是闭集.
\end{proof}

\begin{corollary}
    设 $X$ 是一个紧致 Hausdorff 空间，$A \subset X$.那么 $A$ 是紧的 $\iff$ $A$ 是闭的.
\end{corollary}

\begin{instance}
    设 $X$ 是一个紧致度量空间，$A \subset X$.那么 $A$ 是紧的 $\iff$ $A$ 是闭的.
\end{instance}

\begin{proposition}
    设 $f: X \to Y$ 是从一个紧致空间 $X$ 到一个 Hausdorff 空间 $Y$ 的双射连续映射.那么 $f$ 是一个同胚.
\end{proposition}
\begin{proof}
    只需证明 $f$ 是一个闭映射.假设 $C$ 是 $X$ 中的一个闭集.
    \[
        \left.
        \begin{array}{r}
            X \text{ 紧致} \\
            C \text{ 闭}
        \end{array}
        \right\} \implies C \text{ 在 } X \text{ 中是紧的} \implies f(C) \text{ 在 } Y \text{ 中是紧的}
    \]
    又因为 $Y$ 是 Hausdorff 空间，所以 $f(C)$ 在 $Y$ 中是闭的.因此 $f$ 是一个闭映射.
\end{proof}

\begin{proposition}
    设 $f: X \to Y$ 是从一个紧致空间 $X$ 到一个 Hausdorff 空间 $Y$ 的单射连续映射.那么 $f$ 是一个嵌入.
\end{proposition}
\begin{proof}
    只需证明 $f: X \to f(X)$ 是一个闭映射.因为 $f(X)$ 是 Hausdorff 空间的子空间，所以它也是 Hausdorff 的.因此我们可以应用与前面相同的论证.
\end{proof}

\begin{instance}
    如果 $f: S^1 \to \mathbb{E}^n$ 是单射且连续的，那么它是一个嵌入.
\end{instance}

\begin{lemma}
    设 $X$ 是一个 Hausdorff 空间，$\{C_\lambda\}_{\lambda \in \Lambda}$ 是 $X$ 中一族紧致子集.那么 $\bigcap_{\lambda \in \Lambda} C_\lambda$ 在 $X$ 中是紧的.
\end{lemma}
\begin{proof}
    \[
        \left.
        \begin{array}{r}
            X \text{ Hausdorff} \\
            C_\lambda \text{ 紧致}
        \end{array}
        \right\} \implies C_\lambda \text{ 闭} \implies \bigcap_{\lambda \in \Lambda} C_\lambda \text{ 闭}
    \]
    固定一个 $\lambda_0 \in \Lambda$.由于 $\bigcap_{\lambda \in \Lambda} C_\lambda$ 是紧致空间 $C_{\lambda_0}$ 的一个闭子集，所以它在 $C_{\lambda_0}$ 中是紧的，因此在 $X$ 中也是紧的.
\end{proof}

\begin{lemma}[Lebesgue 引理]
    设 $(X, d)$ 是一个紧致度量空间，$\mathcal{O}$ 是 $X$ 的一个开覆盖.那么存在一个 $\varepsilon > 0$ 使得对任意 $x \in X$，都存在一个 $U \in \mathcal{O}$ 满足 $B(x, \varepsilon) \subset U$.
\end{lemma}
\begin{remark}
    这个 $\varepsilon$ 依赖于 $(X, d)$ 和 $\mathcal{O}$，但不依赖于 $x \in X$.这个 $\varepsilon$ 称为\textbf{Lebesgue 数}.
\end{remark}
\begin{proof}
    如果 $X \in \mathcal{O}$，那么引理显然成立.
    假设 $X \notin \mathcal{O}$.因为 $X$ 是紧的，所以 $\mathcal{O}$ 有一个有限子覆盖 $\{U_1, U_2, \dots, U_n\}$.
    令 $C_i = X - U_i$，则 $C_i$ 是闭集.
    定义函数 $f_i: X \to \mathbb{E}^1$ 为 $f_i(x) = d(x, C_i) = \inf\{d(x, y) \mid y \in C_i\}$.那么 $f_i$ 是连续的.
    令 $f: X \to \mathbb{E}^1$ 为 $f(x) = \frac{1}{n} \sum_{i=1}^n f_i(x)$.那么 $f$ 也是连续的.
    \paragraph{断言 1：}\uline{对任意 $x \in X$，$f(x) > 0$}.
    对任意 $x \in X$，由于 $\{U_i\}$ 是一个覆盖，所以 $x$ 至少属于一个 $U_i$.假设 $x \in U_i$.因为 $U_i$ 是开集，所以存在 $\delta > 0$ 使得 $B(x, \delta) \subset U_i$，这意味着 $d(x, C_i) \ge \delta$，即 $f_i(x) \ge \delta$.
    因此 $f(x) \ge \frac{1}{n} f_i(x) \ge \frac{\delta}{n} > 0$.
    \paragraph{断言 2：}\uline{$f: X \to \mathbb{E}^1$ 取得一个最小值 $\varepsilon > 0$.}
    因为 $X$ 是紧的，所以连续函数 $f$ 在 $X$ 上取得最小值.根据断言 1，这个最小值 $\varepsilon$ 必须大于 0.
    \paragraph{断言 3：}\uline{$\forall\, x\in X,\exists\, U_i\in\{U_1, U_2, \dots, U_n\}$ 使得 $B(x, \varepsilon) \subset U_i$.}\par
    如若不然，则$\exists\, x_0\in X$，使得任意$U_i$都不包含$B(x_0,\varepsilon)$，则$d(x_0,c_i)<\varepsilon,\ie f_{c_i}(x_0)<\varepsilon$.\par
    因此$f(x_0)=\frac{1}{n}\sum_{i=1}^n f_{c_i}(x_0)<\frac{1}{n}\cdot(n\varepsilon)=\varepsilon$.
    \paragraph{证明引理}
    对任意 $x \in X$，我们有 $f(x) \ge \varepsilon$.这意味着 $\frac{1}{n} \sum_{i=1}^n f_i(x) \ge \varepsilon$，所以 $\sum_{i=1}^n f_i(x) \ge n\varepsilon$.
    根据均值原理，必然存在某个 $j$ 使得 $f_j(x) \ge \varepsilon$，即 $d(x, C_j) \ge \varepsilon$.
    这意味着 $B(x, \varepsilon) \cap C_j = \emptyset$，所以 $B(x, \varepsilon) \subset X - C_j = U_j$.引理得证.
\end{proof}

\begin{instance}
    $[0,1]\subset \mb{E}^1$是一个紧致度量空间，$\mcal{O}$是一个$[0,1]$的开覆盖，则$\exists\,n\in\mb{Z}^+,\st \left[\frac{i}{n},\frac{i+1}{n}\right]$被包含在某些$U\in\mcal{O}$中，$i=0,1,\cdots,n-1$.
\end{instance}




\section{紧致空间的乘积}

\begin{lemma}[管道引理]
    设 $X, Y$ 是拓扑空间.那么 $X \times Y$ 是紧的 $\iff$ $X$ 和 $Y$ 都是紧的.
\end{lemma}
\begin{proof}
    “$\Rightarrow$” 假设 $X \times Y$ 是紧的.投影映射 $\pi_1: X \times Y \to X$ 和 $\pi_2: X \times Y \to Y$ 都是连续满射.因为紧致空间的连续像是紧的，所以 $X = \pi_1(X \times Y)$ 和 $Y = \pi_2(X \times Y)$ 都是紧的.

    “$\Leftarrow$” (Tube Lemma 的证明思路) 假设 $X$ 和 $Y$ 都是紧的.令 $\mathcal{O}$ 是 $X \times Y$ 的一个开覆盖.
    固定一点 $x \in X$.子空间 $\{x\} \times Y$ 与 $Y$ 同胚，因此是紧的.由于 $\mathcal{O}$ 也覆盖 $\{x\} \times Y$，所以存在 $\mathcal{O}$ 的一个有限子集，记作 $\{W_1, \dots, W_k\}$，它覆盖了 $\{x\} \times Y$.令 $W_x = \bigcup_{i=1}^k W_i$，则 $W_x$ 是包含 $\{x\} \times Y$ 的开集.

    对每个 $y \in Y$，点 $(x, y)$ 属于某个 $W_i$.因为 $W_i$ 是乘积拓扑中的开集，所以存在 $x$ 的邻域 $U_y$ 和 $y$ 的邻域 $V_y$ 使得 $(x, y) \in U_y \times V_y \subset W_i$.
    集合 $\{V_y \mid y \in Y\}$ 是紧致空间 $Y$ 的一个开覆盖，因此它有一个有限子覆盖 $\{V_{y_1}, \dots, V_{y_n}\}$.
    令 $R_x = \bigcap_{i=1}^n U_{y_i}$.这是一个包含 $x$ 的开集（有限个开集之交）.我们称 $R_x \times Y$ 为一个“管道”(tube).
    我们有
    \[
        R_x \times Y = R_x \times \left(\bigcup_{j=1}^n V_{y_j}\right) = \bigcup_{j=1}^n (R_x \times V_{y_j})
    \]
    由于 $R_x \subset U_{y_j}$，我们有 $R_x \times V_{y_j} \subset U_{y_j} \times V_{y_j}$.而根据构造，$U_{y_j} \times V_{y_j}$ 包含于 $W_x$ 的某一个构成元素中.因此 $R_x \times Y \subset W_x$.这意味着每个“管道” $R_x \times Y$ 都可以被 $\mathcal{O}$ 的有限个元素所覆盖.

    现在，集合 $\{R_x \mid x \in X\}$ 是紧致空间 $X$ 的一个开覆盖.因此，它有一个有限子覆盖 $\{R_{x_1}, \dots, R_{x_m}\}$.
    于是，
    \[
        X \times Y = \left(\bigcup_{i=1}^m R_{x_i}\right) \times Y = \bigcup_{i=1}^m (R_{x_i} \times Y)
    \]
    我们已经知道每个 $R_{x_i} \times Y$ 都可以被 $\mathcal{O}$ 的一个有限子集（即构造 $W_{x_i}$ 的那些开集）所覆盖.因此，$X \times Y$ 可以被这些有限子集的并集所覆盖，这仍然是一个 $\mathcal{O}$ 的有限子集.
    所以 $X \times Y$ 是紧的.
\end{proof}

\begin{corollary}[Tychonoff 定理]
    设 $X_1, \dots, X_n$ 是拓扑空间.那么乘积空间 $X_1 \times \dots \times X_n$ 是紧的 $\iff$ 每个 $X_i$ 都是紧的.
\end{corollary}
\begin{proof}
    可以通过对 $n$ 进行数学归纳法来证明.
\end{proof}



